\documentclass[11pt]{article}
\usepackage{amsmath, amssymb, amsthm}
\usepackage[margin=1in]{geometry}
\usepackage{hyperref}

\newcommand{\Z}{\mathbb{Z}}
\newcommand{\Sym}{S}
\newcommand{\RC}{\mathrm{RC}}
\newcommand{\id}{\mathrm{id}}
\newcommand{\inv}{\mathrm{inv}}
\newcommand{\des}{\mathrm{des}}
\newcommand{\code}{\mathrm{code}}
\newcommand{\Sch}{\mathfrak{S}}
\newcommand{\tens}{\otimes}

\theoremstyle{definition}
\newtheorem{definition}{Definition}
\newtheorem{proposition}{Proposition}
\newtheorem{remark}{Remark}

\title{Specification: Bounded RC Factor Algebra}
\author{}
\date{}

\begin{document}
\maketitle

\section{Overview}

The \textbf{Bounded RC Factor Algebra} of size $n$ is a $\Z$-algebra whose additive basis consists of tensor monomials in full Grassmannian RC graphs of prescribed row counts.  It provides a combinatorial framework for multiplying Schubert polynomials via squash products and normalized tensor keys.  The construction builds on the Molev--Sagan type formula for double Schubert polynomials~\cite{samuelmolev}.

\section{Additive Basis}

\begin{definition}[Full Grassmannian RC graph]
An RC graph $G$ with $k$ rows is \emph{full Grassmannian} if either $\pi(G) = \id$ or $\des(\pi(G)) = \{k-1\}$ (i.e.\ the associated permutation has at most one descent, occurring at position $k-1$).
\end{definition}

\begin{definition}[Basis keys]
Fix an integer $n \geq 1$.  A \textbf{basis key} of size $n$ is a tensor
\[
  E_1 \tens E_2 \tens \cdots \tens E_{n-1} \tens G_n
\]
where:
\begin{itemize}
  \item For $1 \leq i \leq n-1$, each $E_i$ is an $i$-row full Grassmannian RC graph whose permutation lies in $\Sym_{i+1}$ (equivalently, $E_i$ is an elementary symmetric RC graph: $\pi(E_i) \in \Sym_{i+1}$ with $\des(\pi(E_i)) \subseteq \{i-1\}$).
  \item $G_n$ is an $n$-row full Grassmannian RC graph (so $\des(\pi(G_n)) \subseteq \{n-1\}$).
\end{itemize}
Any factor with $\pi = \id$ may be omitted from the tensor (the identity RC graph acts as a unit under the squash product).

The algebra is the free $\Z$-module on all such keys (with fixed size $n$), modulo the normalization relations described below.
\end{definition}

\section{Product}

\subsection{Squash product on RC graphs}

Given RC graphs $A$ (with $m$ rows) and $B$ (with $k$ rows, $k \geq m$), the \textbf{squash product} $A \cdot_{\mathrm{sq}} B$ is computed by:
\begin{enumerate}
  \item Form the \emph{disjoint union}: place $B$ next to $A$ by shifting its column indices into a sufficiently high range that the corresponding Schubert polynomials commute (i.e.\ the permutations act on disjoint positions).  The row counts are unchanged; the number of rows of the result increases only to accommodate the combined code length.
  \item Repeatedly \emph{zero out the last row} (applying the transition/exchange property to absorb the bottom row into the upper rows) until the result has $k$ rows.
\end{enumerate}
The result is a $k$-row RC graph whose permutation satisfies $\pi(A \cdot_{\mathrm{sq}} B) = \pi(A) \cdot \pi(B)$ in the appropriate sense.

\subsection{Key-to-RC-graph evaluation}

A key $K = (F_1, F_2, \ldots, F_m)$ of size $n$ evaluates to a single RC graph by left-to-right squash product:
\[
  \RC(K) = ((\cdots((F_1 \cdot_{\mathrm{sq}} F_2) \cdot_{\mathrm{sq}} F_3) \cdots) \cdot_{\mathrm{sq}} F_m),
\]
resized to $n$ rows.  Each intermediate accumulator is resized upward as needed before squashing.

\subsection{Multiplication of keys}

Given keys $K_L$ and $K_R$ of the same size $n$, their product key is:
\[
  K_L \cdot K_R = \mathrm{normalize}(K_L, K_R),
\]
where the concatenation $(F_1^L, \ldots, F_a^L, F_1^R, \ldots, F_b^R)$ is passed through the normalization procedure.

\subsection{Normalization}

The normalization loop transforms an arbitrary tuple of full Grassmannian RC graphs into the canonical form $E_1 \tens \cdots \tens E_{n-1} \tens G_n$.  It repeats until stable:

\begin{enumerate}
  \item \textbf{Sort and merge.}  Sort factors by row count (ascending).  Merge adjacent factors of the same row count via squash product.  If a merged result is the identity, remove it.

  \item \textbf{Strip identities.}  Remove any factor with $\pi = \id$ and normalize each remaining RC graph by snapping the number of rows to the length of the Lehmer code.

  \item \textbf{Peel oversized factors.}  Scan factors from right to left.  If factor $F_i$ has $\ell(F_i) < n$ rows but $\pi(F_i) \notin \Sym_{\ell(F_i)+1}$, then:
  \begin{enumerate}
    \item Resize $F_i$ to $\ell(F_i) + 1$ rows.
    \item Decompose it as (base, grass) via the squash decomposition, where the Grassmannian part has a single descent at the new bottom row.
    \item Repeat the squash decomposition on the base as long as it is not yet full Grassmannian.
    \item Replace $F_i$ with the resulting sequence of smaller full Grassmannian factors.
    \item Restart the loop from step~1.
  \end{enumerate}
\end{enumerate}

The loop terminates when no further peeling is needed, producing a sorted tuple of full Grassmannian factors with non-decreasing row counts.

\subsection{Multiplication of elements}

Given elements $a = \sum_i \alpha_i K_i$ and $b = \sum_j \beta_j K_j$, the product is
\[
  a \cdot b = \sum_{i,j} \alpha_i \beta_j \, (K_i \cdot K_j),
\]
where each $K_i \cdot K_j$ is computed via key concatenation and normalization as above.

\section{Elementary Symmetric Elements}

\begin{definition}
The elementary symmetric element $e_p^{(k)}$ (in size $n$) is defined as
\[
  e_p^{(k)} = \sum_{G} [G],
\]
where the sum is over all $k$-row RC graphs $G$ with $\pi(G) = s_{k-p+1} s_{k-p+2} \cdots s_k$ (i.e.\ $\code(\pi(G)) = (\underbrace{0,\ldots,0}_{k-p},\underbrace{1,\ldots,1}_{p})$), and $[G]$ denotes the basis element for the key $(G)$ of size $n$.
\end{definition}

The SEM (standard elementary monomial) representation of any polynomial in the first $n$ variables can be written as a linear combination of products $e_p^{(k)}$ for $1 \leq k \leq n$, $1 \leq p \leq k$ with distinct $k$.

\section{Schubert Elements}

The representation of a Schubert polynomial $\Sch_w$ in the bounded RC factor algebra depends on whether $w \in \Sym_n$ or $w$ merely has Lehmer code of length $\leq n$.

\subsection{Case 1: $w \in \Sym_n$ (code length $\leq n-1$)}

When $w \in \Sym_n$, the Schubert element has \textbf{no Grassmannian part} ($G_n = \id$).  It is defined purely via the SEM expansion:

\[
  \Sch_w = \sum_{\mathbf{a}} c_{\mathbf{a}} \, e_{a_1}^{(1)} \cdot e_{a_2}^{(2)} \cdots e_{a_{n-1}}^{(n-1)},
\]
where the sum is the SEM representation of $\Sch_w$ and the $c_{\mathbf{a}}$ are integer coefficients.

\subsection[Case 2: code length = n]{Case 2: code length $= n$ (not in $\Sym_n$, but code fits in $n$ rows)}

When $w \notin \Sym_n$ but the Lehmer code of $w$ has length $\leq n$, the Schubert element uses the \textbf{Grassmannian part} and is computed via the full CEM factorization.

The key observation is that the $n$-variable polynomial ring $\Z[x_1, \ldots, x_n]$ is a free module over $\Lambda_n = \Z[e_1^{(n)}, e_2^{(n)}, \ldots, e_n^{(n)}]$ (the ring of symmetric polynomials in $n$ variables), with basis $\{\Sch_w : w \in \Sym_n\}$.  Concretely, for any $w$ with code length $\leq n$:

\[
  \Sch_w = \sum_{\substack{v \in \Sym_n \\ G \text{ Grass.}}} c_{v,G} \; \Sch_v \cdot [G],
\]
where:
\begin{itemize}
  \item $v$ ranges over elements of $\Sym_n$,
  \item $G$ ranges over $n$-row Grassmannian RC graphs,
  \item each $\Sch_v$ is expanded in the $\Sym_n$ CEM basis (Case~1 above),
  \item $[G]$ is the basis element for the single-factor key $(G)$ of size $n$,
  \item the coefficients $c_{v,G}$ come from the full CEM decomposition.
\end{itemize}

Concretely, one computes the CEM representation of $\Sch_w$ and splits each monomial's factors into those with $< n$ variables (the $\Sym_n$ part) and those with $n$ variables (the Grassmannian/symmetric part).  The $\Sym_n$ part is expanded via Case~1, and the Grassmannian part is expanded using the full CEM decomposition with the appropriate partition.  The results are multiplied and summed.

\section{Coproduct}

The coproduct on a key $K$ of size $n$ is defined by vertical cuts:
\[
  \Delta(K) = \sum_{i=0}^{n} K^{\leq i} \tens K^{> i},
\]
where for each factor $F_j$ of $K$, we resize $F_j$ to $n$ rows and perform a vertical cut at column $i$, producing a left factor (columns $\leq i$) and a right factor (columns $> i$), each normalized.

\section{Key-to-RC-Graph Correspondence}

Each basis key evaluates to a unique RC graph via iterated squash product (\S3.2).  Different keys may evaluate to the same RC graph; the normalization procedure ensures each key is canonical.

The map $K \mapsto \RC(K)$ extends linearly:
\[
  \sum_i \alpha_i K_i \;\mapsto\; \sum_i \alpha_i \, \RC(K_i).
\]
This allows verification: the product in the bounded RC factor algebra, when mapped to RC graphs, should reproduce the standard Schubert polynomial multiplication.

\bibliographystyle{plain}
\bibliography{../molev-sagan-paper}

\end{document}
